% !TEX root = AImath.v5.tex
% !TEX root = draft.tex
\chapter*{ 结语\,\, 旅程的尾声}
 \begin{introduction}
  \item 数学与人工智能的深度融合回顾
  \item 从理论到实践的智能发展历程
  \item 未来智能研究的方向与挑战
  \item 数学思维在智能时代的永恒价值
  \item 人类智慧与机器智能的和谐共生


\end{introduction}
%	通过本次学习，我们小组对计算机视觉以及计算机图形学有了更加深入的了解，也对深度学习和神经网络模型有了一个初步的“不打不相识”。我们从文艺复兴时代的代数学起步，从基本的线性代数知识一步步跨越到了今日具有巨大势头的人工智能和深度学习，对期间科技与知识的进步感到赞叹不已。
%	我们对计算机科学的几个方面有了一定的认识，作为软件工程学院的学生们，我们也越来越期待以后可以接触到更多的计算机技术，以用来解决实际的问题。作为国家未来的工程学人才，我们会时刻保持对新鲜技术的渴望，永远坚持让明天的科技将于今日到来，以造福世界上的每一个人类。
%	很感谢您阅读此篇文章。
这场数学与人工智能携手之旅迎来了尾声。
%
%“诺贝尔物理学奖获得者狄拉克也说“数学家参与游戏并在游戏中创造规律”；而在物理学家参与的游戏中，“规律由大自然提供”。只是“随着时间的流逝，我们会愈加明显地察觉到，那些在数学家眼里有趣的规律，同样也是大自然的选择”。
%
%海因里希王子在少女的笔记本里瞥见了“奇妙的象形文字”，“巫术柱”上成列的“占卜符号”和“咒语”。大自然的规律似乎隐藏在这些符号与咒语中，那些“巫术”符号也被解读为现代数学语言中的三角形和圆形。而伽利略认为，大自然就是用数学语言写成的。复数就是个例子。一位富于幻想的数学家在16世纪的一个宿命时刻创造了复数---一种没有意义的符号，而复数就是现代量子力学中“微粒运行的场所”。”

%
%“可那个想象中的“甚至在空气之外的”游戏有着无法否认的纯美。这种美，源自鲍耶、罗巴切夫斯基、高斯和黎曼这些天才在宿命时刻所创造的理论。他们革新了我们对宇宙几何的看法和时空观，直到我们在相对论中找到了表达世界的方式。狄拉克认为，相对论“在描述大自然时，展现了前所未有的数学之美”。正是它“令人惊叹的数学之美”使物理学家们接受了爱因斯坦的理论，同时也为物理学带来了美学上的深刻改变”
%
%“绘画艺术之美源于自由的创造力，数学之美亦是如此。19世纪德国数学家格奥尔格·康托尔是现代数学无穷理论的创造者，他也曾说过，“数学的本质就在于它的自由”。除此之外，我们通过这本书还想传达什么呢？当远古先祖在骨头上刻痕，数字出现在他们指间，也就是当数字开始阐释世界的时候，我们看见，在数学发明的高光下闪耀的宿命时刻，在人类历史进程中划分先后的那些瞬间，自由在其中涌现。”
%
%%视频网站的魔法
%%
%%我逐渐沉迷于在网上输出大量知识普及内容的先驱者。全靠他们以及其他知识普及者，我才接触到新可乐的故事、阿施从众实验、米尔格兰姆服从性实验、津巴多的斯坦福监狱实验、安慰剂与反安慰剂效应，以及利贝和海恩斯有关自由意志的实验。我自己也开了个博客，写了一篇博文来讲述这些实验[13]。
%%这也正是科普的美妙之处，科普总是让人想知道得更多。我非常感兴趣，阅读了所有这些伟大的研究者的著作。我发现了一个全新的宇宙，其中探寻了我们为什么会有某种念头，以及怎样才能更好地思考我们心中的想法。我学到了很多，但尤其重要的是，我对自身无知的广度把握得越来越准确[14]。
%
%
%显然，任何作者都没有办法做到面面俱到，不可能讨论到所有关键要点、介绍到所有重要人物，或涉及所有亟待解决的数学问题。作者每次都必须做出选择，而这些选择又要受到内在一致性、题材的复杂程度、作者的兴趣和专业知识的限制，还要受到完全人为的字母顺序的限制。这类书的选题策划方案决定了它难免挂一漏万，而大量的好素材最终都不得不忍痛割爱了。
%
%这样一来，本书就成为一个人只身面对浩瀚数学宇宙的感悟。跟随本书在数学知识的海洋中遨游，只能经历无数条路径中的一条，而且我也认为我所选择的从A到Z的顺序并不是最完美的路径。
%
%数学，有趣而美妙。抛开限制不谈，我仍然希望本书至少能够展示数学这门魅力无穷的学科的概貌。正如19世纪数学家索菲亚·柯瓦列夫斯卡娅所说：​“许多无缘更深入认识数学的人士把数学与算术混为一谈，而且还误认为它是一门枯燥无味的科学。然而实际上，它是一门需要最强大想象力的科学。​”[1]也许这本书能够再现5世纪希腊哲学家普罗克洛斯（Proclus）的高尚情怀：​“单凭数学便能重振生机，唤醒灵魂……赋予其生命，化想象为现实，变黑暗为智慧的光芒。​”[2]
%
%
%旅途仍在继续
%但这些形形色色的知识缺少整体结构。庞加莱曾这样说过：“事实的堆砌与一门科学的距离并不比一堆石头与一座房屋的距离近。”
%
%%我开始寻求一个关于理论的理论，某种能够让我将形形色色的知识结合起来并得到更好的理解的东西。在很长一段时间中，我并没有意识到答案就在一个公式之中，尽管我一直都在研究它;答案就在一个术语“贝叶斯”之中，这个词在我自己的博士论文题目的开头就出现了[15]。2016年初，我逐渐意识到这个公式的重要性。在“骑驴找驴”了很长时间之后，我终于走上了贝叶斯主义的道路。
%
%%没有良知的知识不过是对灵魂的戕害。---弗朗索瓦·拉伯雷(1483或1494—1553)
%
%%我只能请你多花时间思索这个惊人的结论。我斗胆更进一步期望这本书能够打破你此前对逻辑、知识以及各个方面认知的看法。我也期望这本书能够帮助你更好地确定科学方法的局限性。我还期望这本书能够帮助你怀疑自己的自信过度，因为你很可能深受其害。我同样期望这本书能让你隐约看到如何更好地学习与获取知识。
%
%人们在处理“现实”世界的问题时，经常认为数学和哲学与之毫不相干。而人们在谈论日常而具体的事物时，也通常觉得这个领域无须数学博士学位就能理解。这就是严重的自信过度。正如约翰·冯·诺伊曼所说：“如果有人不相信数学是简单的，那是因为他没有意识到人生有多复杂。”我们无法理解乌鸦悖论，无法理解为什么要选择公立医院而不是私人诊所，无法理解指数增长有多疯狂，这种能力的缺失理应迫使我们怀疑在“真实”世界中自认为已经理解的东西。
%
%%但这并不是本书的主要目的。正如在第1章中所说的那样，我近几年对认识论的长久思考让我放弃了所谓的科学方法和频率主义。
%
%%但最重要的是，我希望你能喜欢这段探索贝叶斯主义的基础和推论必经的旅程，并享受到探索众多学科的乐趣。这些学科对贝叶斯主义的理解和诠释来说非常有用，无论是理论计算机还是认知科学，或者演化生物学和统计物理学。我希望你喜欢奥卡姆剃刀的证明、对归纳问题的抽丝剥茧，还有对实在论的质疑。我也希望阅读这本书对你来说是一次不同寻常的旅程，甚至为你打开了新世界的大门，让你留下不可磨灭的回忆。
%
%最重要的是，我希望你能沉浸在着迷和疑问之中，这就是本书的首要目的。
面对未来的不确定性，我们通常力图高瞻远瞩、殚精竭虑，利用尽量多的数据和信息，以实现对未来的精准预测。
但是，真正有效的方法恰恰与之相反，我们更应回顾过去，以拓宽视野。与其执着于把握未来的微小变化，不如研究那些过去避无可避的重大事件。
10年前, 我决心少看点儿预测书，多读些历史书。这是我人生中最明 智的决定。有意思的是，历史读得越多，我对未来就越坦然。当开始 聚焦于永恒不变的常识时，你就不会再痴迷于预测未知，而是花更多 时间去研究那些恒久不变的人类行为。希望本书能给你启发，引导你 朝这个方向前进。
我尽量不给自己不了解的人提供建议，因为每个人都是不同的个体，放之四海而皆准的建议少之又少。

“纸上得来终觉浅，绝知此事要躬行。”


我们以读者熟悉的案例说明方法的边界与取舍，避免抽象堆砌。

\newpage
\nocite{*}
%\printbibliography[heading=bibintoc, title=\ebibname]

